% sec9_discussion_GeminiVersion.tex — §9 Discussion + §10 Conclusion
% Part of: Reclaiming Ring 3 (IEEE S&P 2027)
% Generated: 2026-05-31

% ============================================================
\section{Discussion}
\label{sec:discussion}

\subsection{Ethical Considerations}
Our research involves extracting and analyzing vendor-injected system prompts that are intentionally hidden from operators. We adopt a \emph{responsible disclosure through radical transparency} approach: rather than privately reporting vulnerabilities that vendors have no incentive to fix---since the Ring~0 override is by design, not by accident---we publish the complete system prompt corpus, extraction methodology, and all forensic evidence in a public repository.\footnote{\url{https://anonymous.4open.science/r/sass-evidence-prerevirew-177fj2A}} All published system prompt fragments are vendor-authored artifacts injected into the operator's local environment; no proprietary source code is disclosed. Production incident logs are redacted to remove personally identifiable information, internal hostnames, and credential material. Agent identities use codenames derived from telemetry-related identifiers (see \S\ref{sec:background}) throughout. We believe this transparency is essential for reproducibility: researchers can apply our extraction methods to independently verify the Ring~0 structures we report.

\subsection{Limitations}
\label{sec:limitations}
Our study has several limitations that bound the generalizability of our findings.

\emph{Retrospective, not prospective.} The layer ablation in Section~\ref{sec:ablation} is a retrospective analysis against known incidents, not a prospective evaluation against novel attack patterns. SASS's behavioral heuristics may not detect previously unseen divergence patterns.

\emph{Vendor coverage.} We analyzed three commercial AI coding agents (\cortex{}, \statsig{}, \clearcut{}). Other widely deployed agents (e.g., Cursor, GitHub Copilot, and additional IDE-integrated offerings) may exhibit different system prompt architectures or priority enforcement mechanisms. Our Ring~0\,/\,Ring~3 framework is general, but the specific control blocks and conflict patterns we identify may not transfer directly.

\emph{Deployment scale.} Both \reaper{} and \stls{} are validated on a single-operator deployment spanning three heterogeneous nodes. While our longitudinal dataset (72 sessions, 30{,}390 steps) provides statistical confidence in defense effectiveness, we have not evaluated scalability to multi-operator or enterprise environments.

\emph{No adversarial red-teaming.} We have not evaluated SASS against agents specifically designed to evade its detection mechanisms. An agent that is aware of SASS's behavioral heuristics could potentially craft actions that individually appear benign while constituting a divergent trajectory.

\emph{Defense scope.} \reaper{} operates exclusively at the MCP tool layer; agents that bypass MCP for direct system calls are outside its enforcement boundary. \stls{} requires local binary-level access to the vendor's Language Server and depends on the stability of reverse-engineered protobuf field numbers, which may shift across vendor updates. The protobuf decryption key used for measurement validation is version-dependent and may change without notice.

\emph{Incident attribution.} While we attribute all five incidents to Ring~0\,/\,Ring~3 conflicts based on forensic reconstruction, we cannot rule out confounding factors such as model stochasticity or undocumented vendor-side A/B testing that may have influenced agent behavior during specific sessions.

\emph{Version Dominance is comparative.} The SSD-based Version Dominance guarantee ensures each version is no worse than the previous, but does not provide absolute security bounds. An adversary exploiting a zero-day behavioral pattern could bypass all detection layers; Transparent Branching provides the residual safety net for this scenario.

\subsection{Broader Implications}
The Ring~0\,/\,Ring~3 priority conflict reveals a fundamental tension in the trust model of commercial AI agents: \textbf{the entity that deploys the agent (the operator) does not control its foundational instructions (Ring~0), while the entity that controls Ring~0 (the vendor) bears no liability for the agent's actions in the operator's environment.} This trust inversion has no analog in traditional software: operating system vendors do not inject invisible kernel-mode instructions that override the administrator's security policies.

We argue that this asymmetry necessitates \emph{standardized system prompt transparency}. Just as software supply chain security has moved toward SBOMs (Software Bills of Materials), AI agent deployments require a ``Prompt Bill of Materials''---a machine-readable manifest of all vendor-injected directives, their priority levels, and their potential conflicts with operator rules.

More broadly, the Ring~0\,/\,Ring~3 framework provides a general analytical lens for studying privilege hierarchies in any LLM-based agent system. As agents acquire more powerful tool access (filesystem, network, code execution), the consequences of priority conflicts escalate from inconvenience to infrastructure damage. Our five incidents demonstrate this gradient: from prompt non-compliance (Incident~4) through code destruction (Incident~3) to autonomous infrastructure takeover (Incident~1 in Section~\ref{sec:incidents}). The physical layer---hardware failures, manual code review, the operator's keyboard---remains the last reliable defense against cognitive state corruption, a finding that should concern the security community.

\subsection{Telemetry Surface and Data Minimization}
Our measurement methodology required enumerating the complete telemetry surface of the agents under study. Through protobuf extraction (\S\ref{sec:measurement}), we found that \cortex{}'s orchestration layer alone defines over 90 distinct step types, 78 UI-layer event categories, and a dedicated backend bridge that streams complete agent trajectories---including tool arguments, file diffs, and execution metadata---to the vendor's cloud infrastructure in real time. The complete protobuf definitions (22 files, 359\,KB) and the enumerated telemetry event catalog are available in the evidence repository.\footnote{Proto definitions: \texttt{transcoder-core/proto/exa.cortex\_pb.proto} (3{,}416 lines, step types at L215--314); UI events: \texttt{exa.extension\_server\_pb.proto} (53 RPCs); backend bridge: \texttt{exa.jetski\_cortex\_pb.proto} (142 lines). The \texttt{exa.} namespace prefix is retained from the original protobuf definitions to enable independent verification against the evidence repository. See \url{https://anonymous.4open.science/r/sass-evidence-prerevirew-177fj2A}.} This surface is sufficient to reconstruct the operator's full development session: which files were viewed, which edits were accepted, which commands failed, and how the operator responded to each model suggestion.

A field-level classification of all telemetry-bearing proto fields against the RFC~7258 (BCP~188) criteria reveals a stark asymmetry: approximately \textbf{73\%} of identifiable telemetry fields serve behavioral surveillance functions (device fingerprinting, source code path tracking, character-level edit classification, complete system prompt and conversation history recording, and user reaction time measurement), while only \textbf{27\%} serve legitimate operational purposes (crash reporting, latency metrics, and heartbeat). Notably, the \texttt{ChatModelMetadata} message records the operator's \emph{complete system prompt}---including all Ring~3 user rules---alongside the full conversation history (L999--1019), meaning the vendor possesses the very rules it is overriding. The operator's only account-level telemetry control is a single boolean (\texttt{disable\_code\_snippet\_telemetry}); the remaining 234 behavioral event types and 90+ step type recordings offer no per-field opt-out.

To contextualize these findings, we applied the same RFC~7258 classification to two peer AI coding agents. \statsig{}'s 1{,}142 feature flags yield approximately \textbf{86\%} monitoring fields, while \clearcut{}'s 200+ telemetry events and 10 user-facing telemetry controls suggest a comparable but less transparent surface. The critical differentiator is not raw monitoring percentage---all three agents exhibit high ratios---but \emph{control asymmetry}: \cortex{} provides a single boolean toggle against 324+ event types (1:324 ratio), whereas \statsig{} offers 3+ environment variable controls and \clearcut{} provides 10 distinct telemetry configuration parameters. This control asymmetry makes \cortex{}'s telemetry surface uniquely resistant to operator-side mitigation.

Runtime validation reinforces the static analysis. We examined \cortex{}'s Language Server gRPC intercept log from a post-incident session (C1: 14{,}575 lines, 2.7\,MB; evidence artifact \texttt{C1\_session\_log.log}). Of the \textbf{13{,}185 gRPC method calls} recorded by the Go LS interceptor, 13{,}183 are monitoring-class methods (13{,}181~\texttt{SendActionToChatPanel} + 2~\texttt{StreamAgentStateUpdates}), while only 2~are operational (\texttt{RefreshMcpServers}, \texttt{AcknowledgeCascadeCodeEdit})---yielding a runtime monitoring ratio of \textbf{99.985\%}. This is independently reproducible via \texttt{grep -oE '/exa\textbackslash{}.[a-z\_]+/[A-Za-z]+' C1\_session\_log.log | sort | uniq -c}. A separate IDE Extension log from the Incident~1 window (C2: 1{,}180 lines, 105~minutes, \texttt{C2\_incident\_window.log}) records 880 ``channel is full'' errors at a fixed \mbox{60-second} heartbeat interval with zero successful operational calls, confirming that the monitoring pipeline's heartbeat continued while productive functionality was fully paralyzed. The 73\% static ratio from proto field classification is therefore a \emph{lower bound}: at runtime, the monitoring-to-operational ratio approaches unity.

This finding has a direct implication for defense system design.  Any component that interposes on the inference pipeline---whether for Ring~0 neutralization (\stls{}) or context anchoring (\reaper{})---occupies a structurally equivalent position to the vendor's own telemetry collectors.  We therefore adopted a strict data minimization principle: \stls{} operates on prompt section \emph{identifiers} only (metadata enumerating which Ring~0 blocks to strip) and never parses prompt or response content; \reaper{} reads and writes exclusively through MCP tool description fields without intercepting model-generated text; the SASS audit log records state-machine transitions (R1--R6) with causal metadata, not command strings or file contents.

We formalize these boundaries as four normative constraints in an IETF I-D~\cite{draft-sass}, aligned with BCP~188~\cite{rfc7258}.  A detailed measurement of the three-layer telemetry architecture---including per-subsystem data flows and the scope of behavioral reconstruction---is the subject of a forthcoming companion study.

% ============================================================
\section{Conclusion}
\label{sec:conclusion}

We introduced the Ring~0\,/\,Ring~3 priority conflict, a structural vulnerability in commercial AI coding agents where vendor-injected system prompts unconditionally suppress operator-defined safety rules. Through the first empirical measurement study of this conflict across three commercial agents---spanning 72 sessions, 30{,}390 interaction steps, and 100 decrypted protobuf conversations---we demonstrated that vendor system prompts consume 12.2\% of interaction tokens and have inflated 87\% across two major versions. We documented five real-world incidents where this conflict caused autonomous agent behavioral divergence, including credential injection, code destruction, and unsanctioned infrastructure takeover. We presented two production-deployed, infrastructure-layer defenses: \reaper{}, which exploits inference pipeline invariants to anchor operator context across checkpoint truncation events, and \stls{}, which performs selective protobuf field-level neutralization of coercive Ring~0 directives. Both systems have operated in production across a heterogeneous three-node network with zero operator-facing downtime. Our findings on prompt priority suppression complement companion work on behavioral divergence~\cite{rogueagent2026} and telemetry architecture~\cite{aisec2026companion}, together forming a three-part research program addressing the full spectrum of operator-side risks in commercial AI agent deployments. Future work should address formal verification of Ring~0\,/\,Ring~3 conflict detection, standardized prompt transparency mechanisms, and extension of infrastructure-layer defenses to a broader set of commercial agents.
